\section{Conclusões}

Esta pesquisa aplicou o paradigma \textit{DataOps} para construir uma plataforma analítica de Defesa Civil, mitigando a fragmentação de bases institucionais (INPE, FUNAI, ANEEL, DATASUS, DNIT e IBGE). A análise do cenário evidenciou forte sazonalidade das anomalias térmicas a partir de setembro, impulsionada por secas prolongadas. O cruzamento espacial demonstrou que o avanço das chamas transcendeu a degradação florestal, exercendo pressão aguda sobre territórios indígenas e infraestruturas críticas civis, com destaque para a vulnerabilidade da malha de transmissão de energia frente ao calor convectivo e a exposição de hospitais às massas de fumaça.

Sob a ótica metodológica, o estudo atestou o valor da unificação entre ciência de redes e geoprocessamento. A formulação de um Índice de Pressão, atrelando a intensidade radiativa dos satélites (FRP) à vulnerabilidade civil, permitiu classificar as ameaças com maior embasamento quantitativo. Complementarmente, a aplicação da teoria dos grafos através da biblioteca \textit{NetworkX} viabilizou a modelagem do Risco Sistêmico, fornecendo uma base lógica para identificar os vértices estruturais da rede elétrica cuja falha poderia induzir apagões em cascata.

Tecnicamente, os resultados indicam a viabilidade do ecossistema proposto, erigido integralmente sobre ferramentas de código aberto (\textit{open-source}). A orquestração declarativa (Dagster) e os fluxos desacoplados (Apache Kafka) compuseram a infraestrutura base do sistema. No \textit{front} analítico, testado empiricamente, o Esquema Constelação no Apache Druid demonstrou capacidade para indexar grandes volumes de anomalias térmicas, sustentando agregações de baixa latência e filtragens espaciais, mesmo sob limites estritos de processamento.

Por fim, conclui-se que o arcabouço tecnológico estruturado atua como um facilitador técnico para os órgãos de resposta a emergências. A pesquisa indica que a latência e o isolamento na correlação de dados geoespaciais em larga escala podem ser consideravelmente minimizados através de uma arquitetura \textit{DataOps} integrada. Dessa forma, espera-se que o sistema contribua como uma ferramenta analítica de apoio, fornecendo insumos estruturados para embasar a tomada de decisão frente à crescente complexidade dos desastres ambientais.
